\documentclass[12pt]{article}
\usepackage{amsmath}
\usepackage{amsfonts}
\usepackage{float}
\usepackage{fancyhdr}
\usepackage{graphicx}
\usepackage{setspace}
\usepackage[colorlinks=true,linkcolor=blue, citecolor=red]{hyperref}
\usepackage{url}
\usepackage[top=.75in, left=.75in, right=.75in, bottom=1in]{geometry}

% For algorithm
\usepackage{algorithm}
\usepackage{algpseudocode}

% ============ CODE ============
\usepackage{listings}
\usepackage{xcolor}
\definecolor{codegreen}{rgb}{0,0.6,0}
\definecolor{codegray}{rgb}{0.5,0.5,0.5}
\definecolor{codepurple}{rgb}{0.58,0,0.82}
\definecolor{backcolour}{rgb}{0.95,0.95,0.92}

% Styling for the code.
\lstdefinestyle{mystyle}{
    backgroundcolor=\color{backcolour},   
    commentstyle=\color{codegreen},
    keywordstyle=\color{magenta},
    numberstyle=\tiny\color{codegray},
    stringstyle=\color{codepurple},
    basicstyle=\ttfamily\footnotesize,
    breakatwhitespace=false,         
    breaklines=true,                 
    captionpos=b,                    
    keepspaces=true,                 
    numbers=left,                    
    numbersep=5pt,                  
    showspaces=false,                
    showstringspaces=false,
    showtabs=false,                  
    tabsize=2
}
\lstset{style=mystyle}

% Table packages and Unicode / Font support
\usepackage{array}
\usepackage{longtable}
\usepackage{enumitem}
\usepackage{booktabs}
\usepackage[utf8]{inputenc}
\DeclareUnicodeCharacter{2264}{\ensuremath{\le}}
\DeclareUnicodeCharacter{2194}{\ensuremath{\leftrightarrow}}
\usepackage[T1]{fontenc}
\usepackage{lmodern}
\graphicspath{{../../../../templates/}{../../prototype/}{figures/}{./}}
\setlength{\emergencystretch}{3em}
\setlist[enumerate]{label=\arabic*.,leftmargin=*}
\setlist[itemize]{leftmargin=*}

% ============ HEADER AND FOOTER ============
\setlength{\headheight}{29.43912pt}

\pagestyle{fancy}
\fancyfoot{}
\fancyfoot[R]{Page \thepage}

\setlength{\parindent}{0pt}
\renewcommand{\baselinestretch}{1.5}

\newcommand{\coursename}{Introduction of Software Engineering}
\newcommand{\reportname}{Studitopia}
\newcommand{\reporttitle}{\href{https://github.com/RoundFloor/Studitopia.git}{Studitopia}}

\newcommand{\studentname}{%
Nguyen Thanh An (24125024)\\
Nguyen Chi Bao (24125025)\\
Tran Nguyen Khanh Binh (24125026)\\
Tran Canh Anh Tuan (24125048)\\
Nguyen Le Quoc Huy (24125100)
}
\newcommand{\teachername}{Ho Tuan Thanh, Nguyen Thi Minh Tuyen}

% Header
\lhead{\reporttitle}
\rhead{
University of Science - VNUHCM\\
\coursename
}

% Footer
\newcommand{\leftfooter}{}
\lfoot{\leftfooter}

% ============ DOCUMENT ============
\begin{document}

\pagenumbering{roman}
\begin{titlepage}
\newcommand{\HRule}{\rule{\linewidth}{0.5mm}}
\centering

\textsc{\LARGE VNU-HCM}\\[0.5cm]
\textsc{\Large University of Science}\\[0.5cm]
\textsc{\large Faculty of Information Technology}\\[0.5cm]
\textsc{Computer Science}\\[0.5cm]

\HRule \\[0.4cm]
{ 
\huge{\bfseries{\reporttitle}}\\[0.5cm]
\large{\bfseries{UC-17 --- Vote on Posts and Comments}}
}\\[0.4cm]
\HRule \\[0.5cm]

\textbf{\large Subject: \coursename}\\[0.5cm]

\begin{minipage}[t]{0.55\textwidth}
\begin{flushleft} \large
\emph{Students:}\\
\studentname
\end{flushleft}
\end{minipage}
~
\begin{minipage}[t]{0.4\textwidth}
\begin{flushright} \large
\emph{Supervising Teachers:} \\
\teachername
\end{flushright}
\end{minipage}\\[1cm]

{\large \today}\\[1cm]

\includegraphics[scale=.20]{img/hcmus-logo.png}\\[1cm] 

\vfill
\end{titlepage}	

\tableofcontents
\pagebreak

\pagenumbering{arabic}
\setcounter{page}{1}

\hypertarget{sec:uc-17-vote-on-posts-and-comments}{}
\addcontentsline{toc}{section}{UC-17 --- Vote on Posts and Comments}
\section*{UC-17 --- Vote on Posts and Comments}
\par\medskip

\textbf{Version:} \texttt{1.0}
\par\medskip

\textbf{Performed by:} [Nguyen Le Quoc Huy] | \textbf{Edited by:} [Nguyen Le Quoc Huy]
\par\medskip

\begingroup
\small
\renewcommand{\baselinestretch}{1.0}\selectfont
\renewcommand{\arraystretch}{1.15}
\setlength{\tabcolsep}{6pt}
\begin{longtable}{|>{\raggedright\arraybackslash}p{0.09\textwidth}|>{\raggedright\arraybackslash}p{0.84\textwidth}|}
\hline
\textbf{Field} & \textbf{Value} \\ \hline
\endhead
Use Case ID & UC-17 \\ \hline
Name & Vote on Posts and Comments \\ \hline
Primary Actor & University Student \\ \hline
Supporting Actors & None \\ \hline
Feature Group & Forum Interaction \\ \hline
Priority & Not explicitly stated \\ \hline
Related User Stories & US-19 \\ \hline
Related Requirements & FR-13.2, FR-13.4, FR-17.4, FR-19.1 to FR-19.5, FR-38.1; NFR-PERF-2, NFR-PERF-3, NFR-PERF-4, NFR-PERF-6, NFR-SEC-1, NFR-SEC-6, NFR-SEC-9, NFR-PRIV-7 \\ \hline
Source Documents & \texttt{proposals/userstory.md} (section 6.2); \texttt{proposals/functional\_requirement.md} (US-19); \texttt{proposals/list\_post.md}; \texttt{proposals/post\_management.md}; \texttt{proposals/create\_post.md}; \texttt{proposals/API.md}; \texttt{proposals/response\_code.md}; \texttt{proposals/non\_functional\_requirement.md}; \texttt{proposals/feature\_flow.md}; \texttt{proposals/screen\_flow.md} \\ \hline
\end{longtable}
\endgroup
\par\medskip

\hypertarget{sec:part-a-use-case-model}{}
\addcontentsline{toc}{subsection}{Part A --- Use Case Model}
\subsection*{Part A --- Use Case Model}
\par\medskip

\hypertarget{sec:a1-focused-use-case-diagram}{}
\addcontentsline{toc}{subsubsection}{A1. Focused Use Case Diagram}
\subsubsection*{A1. Focused Use Case Diagram}
\par\medskip

\begin{center}
\includegraphics[width=0.85\textwidth,height=0.75\textheight,keepaspectratio]{"figures/UC-17-vote-posts-comments_1.png"}
\end{center}
\par\medskip

\textbf{Diagram notation.} Person shape = human actor; cloud = supporting system; blue = selected use case; solid line = association; labelled solid arrow = generalization; dashed arrow = include or extend.
\par\medskip

\hypertarget{sec:a2-actor-register}{}
\addcontentsline{toc}{subsubsection}{A2. Actor Register}
\subsubsection*{A2. Actor Register}
\par\medskip

\begingroup
\small
\renewcommand{\baselinestretch}{1.0}\selectfont
\renewcommand{\arraystretch}{1.15}
\setlength{\tabcolsep}{6pt}
\begin{longtable}{|>{\raggedright\arraybackslash}p{0.11\textwidth}|>{\raggedright\arraybackslash}p{0.11\textwidth}|>{\raggedright\arraybackslash}p{0.70\textwidth}|}
\hline
\textbf{Actor} & \textbf{Type} & \textbf{Description} \\ \hline
\endhead
University Student & Primary human actor & A student who upvotes or downvotes a post, comment, or answer so that useful academic content is easier to discover. \\ \hline
\end{longtable}
\endgroup
\par\medskip

\hypertarget{sec:a3-relationships}{}
\addcontentsline{toc}{subsubsection}{A3. Relationships}
\subsubsection*{A3. Relationships}
\par\medskip

\hypertarget{sec:a3-1-associations}{}
\paragraph*{A3.1 Associations}
\par\medskip

\begingroup
\small
\renewcommand{\baselinestretch}{1.0}\selectfont
\renewcommand{\arraystretch}{1.15}
\setlength{\tabcolsep}{6pt}
\begin{longtable}{|>{\raggedright\arraybackslash}p{0.12\textwidth}|>{\raggedright\arraybackslash}p{0.23\textwidth}|>{\raggedright\arraybackslash}p{0.10\textwidth}|>{\raggedright\arraybackslash}p{0.46\textwidth}|}
\hline
\textbf{Actor} & \textbf{Use Case} & \textbf{Type} & \textbf{Description} \\ \hline
\endhead
University Student & UC-17 Vote on Posts and Comments & Primary & Casts, changes, or removes a vote on an eligible post or comment. \\ \hline
\end{longtable}
\endgroup
\par\medskip

\hypertarget{sec:a3-2-generalization-and-specialization}{}
\paragraph*{A3.2 Generalization and Specialization}
\par\medskip

\begingroup
\small
\renewcommand{\baselinestretch}{1.0}\selectfont
\renewcommand{\arraystretch}{1.15}
\setlength{\tabcolsep}{6pt}
\begin{longtable}{|>{\raggedright\arraybackslash}p{0.09\textwidth}|>{\raggedright\arraybackslash}p{0.09\textwidth}|>{\raggedright\arraybackslash}p{0.09\textwidth}|>{\raggedright\arraybackslash}p{0.65\textwidth}|}
\hline
\textbf{Child} & \textbf{Parent} & \textbf{Type} & \textbf{Inherited Behavior or Association} \\ \hline
\endhead
None & None & None & No actor generalization or specialization is named for US-19 in the specified proposals. \\ \hline
\end{longtable}
\endgroup
\par\medskip

\hypertarget{sec:a3-3-include}{}
\paragraph*{A3.3 Include}
\par\medskip

\begingroup
\small
\renewcommand{\baselinestretch}{1.0}\selectfont
\renewcommand{\arraystretch}{1.15}
\setlength{\tabcolsep}{6pt}
\begin{longtable}{|>{\raggedright\arraybackslash}p{0.10\textwidth}|>{\raggedright\arraybackslash}p{0.11\textwidth}|>{\raggedright\arraybackslash}p{0.71\textwidth}|}
\hline
\textbf{Base Use Case} & \textbf{Included Use Case} & \textbf{Meaning} \\ \hline
\endhead
None & None & No distinct included use case is named for US-19 in the specified proposals; no new I-xx ID is introduced. \\ \hline
\end{longtable}
\endgroup
\par\medskip

\hypertarget{sec:a3-4-extend}{}
\paragraph*{A3.4 Extend}
\par\medskip

\begingroup
\small
\renewcommand{\baselinestretch}{1.0}\selectfont
\renewcommand{\arraystretch}{1.15}
\setlength{\tabcolsep}{6pt}
\begin{longtable}{|>{\raggedright\arraybackslash}p{0.10\textwidth}|>{\raggedright\arraybackslash}p{0.09\textwidth}|>{\raggedright\arraybackslash}p{0.09\textwidth}|>{\raggedright\arraybackslash}p{0.62\textwidth}|}
\hline
\textbf{Extending Use Case} & \textbf{Base Use Case} & \textbf{Condition} & \textbf{Added Behavior} \\ \hline
\endhead
None & None & None & No distinct extending use case is named for US-19 in the specified proposals; no new E-xx ID is introduced. \\ \hline
\end{longtable}
\endgroup
\par\medskip

\hypertarget{sec:a4-model-notes}{}
\addcontentsline{toc}{subsubsection}{A4. Model Notes}
\subsubsection*{A4. Model Notes}
\par\medskip

\begin{enumerate}
\item UC-17 covers both target types named in \texttt{API.md}: \texttt{post} and \texttt{comment}.
\item Comments and answers share the \texttt{comment} target type in the API; no separate answer vote endpoint is named.
\item Upvote, downvote, change, and remove are values of one vote action, not separate use cases.
\item Notification of vote activity is a related cross-cutting behavior in \texttt{FR-38.1}, not a source-named included use case; priority and a dedicated vote UI route are not specified in proposals.
\end{enumerate}
\par\medskip

\hypertarget{sec:part-b-use-case-specification}{}
\addcontentsline{toc}{subsection}{Part B --- Use Case Specification}
\subsection*{Part B --- Use Case Specification}
\par\medskip

\hypertarget{sec:b1-brief-description}{}
\addcontentsline{toc}{subsubsection}{B1. Brief Description}
\subsubsection*{B1. Brief Description}
\par\medskip

A university student votes on a post or on a comment or answer so that useful academic content is easier to discover. The student can upvote, downvote, change a vote, or remove a vote, with at most one active vote per target. The recorded vote updates the vote score and ranking, and vote controls appear where the user has permission to view and interact with the target.
\par\medskip

\hypertarget{sec:b2-preconditions}{}
\addcontentsline{toc}{subsubsection}{B2. Preconditions}
\subsubsection*{B2. Preconditions}
\par\medskip

\begin{enumerate}
\item The student is authenticated with an active session.
\item The student can view the selected post or comment.
\item The target is a \texttt{post} or \texttt{comment} and still exists.
\item The student is permitted to vote under privacy, moderation, and authorization rules.
\end{enumerate}
\par\medskip

\hypertarget{sec:b3-postconditions}{}
\addcontentsline{toc}{subsubsection}{B3. Postconditions}
\subsubsection*{B3. Postconditions}
\par\medskip

\begin{enumerate}
\item \textbf{Success:} The vote is recorded as \texttt{1}, \texttt{0}, or \texttt{-1} for the target and requesting user.
\item \textbf{Success:} The vote score and ranking are updated after the vote is recorded.
\item \textbf{Success:} The user's vote state is reflected in the forum view.
\item \textbf{Failure:} The vote is not recorded, the previous vote state remains authoritative, and the student receives a clear error.
\end{enumerate}
\par\medskip

\hypertarget{sec:b4-main-success-scenario-basic-flow}{}
\addcontentsline{toc}{subsubsection}{B4. Main Success Scenario (Basic Flow)}
\subsubsection*{B4. Main Success Scenario (Basic Flow)}
\par\medskip

\begin{enumerate}
\item The student opens a forum post list or post detail page.
\item The system displays the target's current vote score and the student's current vote state.
\item The student selects upvote or downvote on a post, comment, or answer.
\item The frontend sends \texttt{target\_type}, \texttt{target\_id}, and the selected \texttt{vote} value.
\item The backend validates the target type, target existence, vote value, permissions, and rate limits.
\item The system records or replaces the student's vote for that target.
\item The system returns the updated \texttt{vote\_score}.
\item The system refreshes the vote score, ranking, and vote controls in the forum view.
\end{enumerate}
\par\medskip

\hypertarget{sec:b5-alternative-flows}{}
\addcontentsline{toc}{subsubsection}{B5. Alternative Flows}
\subsubsection*{B5. Alternative Flows}
\par\medskip

Unless otherwise stated, an alternative flow resumes at the next step of the basic flow.
\par\medskip

\hypertarget{sec:a1-vote-on-a-comment-or-answer}{}
\paragraph*{A1. Vote on a comment or answer}
\par\medskip

At basic flow step 3, the student votes on a comment or answer. The request uses \texttt{target\_type: 'comment'} with the comment or answer identifier.
\par\medskip

\hypertarget{sec:a2-change-an-existing-vote}{}
\paragraph*{A2. Change an existing vote}
\par\medskip

At basic flow step 3, the student selects the opposite vote direction. The existing vote is replaced for the same target, preserving one active vote per user and target.
\par\medskip

\hypertarget{sec:a3-remove-a-vote}{}
\paragraph*{A3. Remove a vote}
\par\medskip

At basic flow step 3, the student removes the current vote by submitting \texttt{vote: 0}. The score is recalculated without an active vote for that target.
\par\medskip

\hypertarget{sec:a4-vote-from-list-or-detail}{}
\paragraph*{A4. Vote from list or detail}
\par\medskip

At basic flow step 1, the student votes from the post list or from post detail where the corresponding control is displayed.
\par\medskip

\hypertarget{sec:b6-exception-flows}{}
\addcontentsline{toc}{subsubsection}{B6. Exception Flows}
\subsubsection*{B6. Exception Flows}
\par\medskip

\hypertarget{sec:e1-invalid-vote-value-or-target-type}{}
\paragraph*{E1. Invalid vote value or target type}
\par\medskip

The system returns \texttt{400} with \texttt{INVALID\_VOTE} (\texttt{48005}) and does not change the score.
\par\medskip

\hypertarget{sec:e2-duplicate-vote-conflict}{}
\paragraph*{E2. Duplicate vote conflict}
\par\medskip

The system returns \texttt{409} with \texttt{DUPLICATED\_VOTE} (\texttt{48006}). The UI reflects the current vote and offers change or removal rather than creating a second active vote.
\par\medskip

\hypertarget{sec:e3-target-not-found}{}
\paragraph*{E3. Target not found}
\par\medskip

The system returns \texttt{404} with \texttt{POST\_NOT\_FOUND} (\texttt{48001}) or \texttt{COMMENT\_NOT\_FOUND} (\texttt{48002}) and refreshes the unavailable view.
\par\medskip

\hypertarget{sec:e4-unauthorized-or-forbidden-vote}{}
\paragraph*{E4. Unauthorized or forbidden vote}
\par\medskip

The system returns \texttt{401}, or \texttt{403} with \texttt{FORBIDDEN} (\texttt{48010}), and hides or disables the action for unauthorized users.
\par\medskip

\hypertarget{sec:e5-rate-limit-reached}{}
\paragraph*{E5. Rate limit reached}
\par\medskip

The system rejects the vote and asks the student to retry later. \texttt{NFR-SEC-9} requires vote rate limiting, but no rate-limit status or error code is specified in proposals.
\par\medskip

\hypertarget{sec:e6-server-or-database-failure}{}
\paragraph*{E6. Server or database failure}
\par\medskip

The system returns \texttt{500} with \texttt{DATABASE\_ERROR} or \texttt{UNKNOWN\_ERROR} (\texttt{45000}), leaves the previous vote state authoritative, and offers retry where supported.
\par\medskip

\hypertarget{sec:b7-business-rules}{}
\addcontentsline{toc}{subsubsection}{B7. Business Rules}
\subsubsection*{B7. Business Rules}
\par\medskip

\begin{enumerate}
\item \texttt{AC-19.1} / \texttt{FR-19.1}: Users can upvote or downvote a post.
\item \texttt{AC-19.2} / \texttt{FR-19.2}: Users can upvote or downvote a comment or answer.
\item \texttt{AC-19.3} / \texttt{FR-19.3}: A user cannot vote more than once on the same post or comment.
\item \texttt{AC-19.4} / \texttt{FR-19.4}: Users can change or remove their vote.
\item \texttt{AC-19.5} / \texttt{FR-19.5}: The vote score and ranking update after a vote is recorded.
\item The API vote value is \texttt{1 | 0 | -1}, where \texttt{0} is removal; a vote must target \texttt{post} or \texttt{comment}.
\item Authentication, server-side authorization, privacy, and moderation rules determine whether vote controls are available (\texttt{NFR-SEC-1}, \texttt{NFR-SEC-6}, \texttt{NFR-PRIV-7}).
\item Voting is rate-limited to reduce spam and abuse (\texttt{NFR-SEC-9}).
\end{enumerate}
\par\medskip

\hypertarget{sec:b8-special-requirements}{}
\addcontentsline{toc}{subsubsection}{B8. Special Requirements}
\subsubsection*{B8. Special Requirements}
\par\medskip

\begingroup
\small
\renewcommand{\baselinestretch}{1.0}\selectfont
\renewcommand{\arraystretch}{1.15}
\setlength{\tabcolsep}{6pt}
\begin{longtable}{|>{\raggedright\arraybackslash}p{0.16\textwidth}|>{\raggedright\arraybackslash}p{0.77\textwidth}|}
\hline
\textbf{ID} & \textbf{Requirement} \\ \hline
\endhead
NFR-PERF-2, NFR-PERF-3 & The post list must load within 2 seconds, and post detail and comments must display within 2 seconds. \\ \hline
NFR-PERF-4 & A new vote must be reflected in the relevant forum view within 2 seconds. \\ \hline
NFR-PERF-6 & Forum notifications, including vote activity, must be delivered within 2 seconds under normal conditions. \\ \hline
NFR-SEC-1 & Authentication is required before accessing protected interactions. \\ \hline
NFR-SEC-6 & Server-side authorization checks must prevent unauthorized votes. \\ \hline
NFR-SEC-9 & Voting must be rate-limited to reduce spam and abuse. \\ \hline
NFR-PRIV-7 & Privacy settings apply consistently to forum posts, comments, member lists, and search results. \\ \hline
\end{longtable}
\endgroup
\par\medskip

\hypertarget{sec:b9-extension-points}{}
\addcontentsline{toc}{subsubsection}{B9. Extension Points}
\subsubsection*{B9. Extension Points}
\par\medskip

\begingroup
\small
\renewcommand{\baselinestretch}{1.0}\selectfont
\renewcommand{\arraystretch}{1.15}
\setlength{\tabcolsep}{6pt}
\begin{longtable}{|>{\raggedright\arraybackslash}p{0.32\textwidth}|>{\raggedright\arraybackslash}p{0.61\textwidth}|}
\hline
\textbf{Extension Point} & \textbf{Location in Flow of Events} \\ \hline
\endhead
Select vote target & Basic flow steps 1-3; alternatives A1 and A4 \\ \hline
Change or remove vote & Basic flow step 3; alternatives A2-A3 \\ \hline
Permission filtering & Basic flow step 2; exception E4 \\ \hline
Duplicate-vote recovery & Basic flow step 6; exception E2 \\ \hline
Rate-limit recovery & Basic flow step 5; exception E5 \\ \hline
\end{longtable}
\endgroup
\par\medskip

\hypertarget{sec:b10-interface-and-data-notes}{}
\addcontentsline{toc}{subsubsection}{B10. Interface and Data Notes}
\subsubsection*{B10. Interface and Data Notes}
\par\medskip

\begin{itemize}
\item \textbf{Screen/flow path:} \texttt{screen\_flow.md} gives \texttt{Forum -$>$ University -$>$ Major -$>$ Course -$>$ Post}, with \texttt{Trending -$>$ Post}. \texttt{feature\_flow.md} gives \texttt{List Post -$>$ Post}, \texttt{Post -$>$ Upvote / Downvote}, and \texttt{Comment -$>$ Vote}.
\item \textbf{UI components:} \texttt{list\_post.md} recommends \texttt{VoteUpButton}, \texttt{VoteDownButton}, \texttt{PostVoteScoreText}, and \texttt{EmptyPostListState}; \texttt{post\_management.md} recommends \texttt{VoteControl} with current score and user vote state.
\item \textbf{Post list and detail fields:} \texttt{GET /forum/posts} returns \texttt{post\_id}, \texttt{title}, \texttt{post\_type}, \texttt{author}, \texttt{vote\_score}, \texttt{comment\_count}, \texttt{solved}, and \texttt{created\_at} with \texttt{200}; failures are \texttt{400}, \texttt{401}, and \texttt{500}. \texttt{GET /forum/posts/:post\_id} returns \texttt{post}, \texttt{comments}, \texttt{true\_response}, \texttt{user\_vote: -1 | 0 | 1}, and \texttt{saved} with \texttt{200}; failures are \texttt{404 POST\_NOT\_FOUND} (\texttt{48001}), \texttt{401}, and \texttt{500 (DATABASE\_ERROR | UNKNOWN\_ERROR)}, with \texttt{FORBIDDEN: 48010} also declared by \texttt{studitopia.get\_forum\_post}.
\item \textbf{API endpoint and failures:} \texttt{POST /forum/votes} with body \texttt{\{ target\_type: 'post' | 'comment', target\_id: string, vote: 1 | 0 | -1 \}} returns \texttt{200} with \texttt{data: \{ vote\_score: number \}}. Failures are \texttt{400} with \texttt{(INVALID\_VOTE)[]} (\texttt{48005}), \texttt{409} with \texttt{(DUPLICATED\_VOTE)[]} (\texttt{48006}), \texttt{404} with \texttt{(POST\_NOT\_FOUND | COMMENT\_NOT\_FOUND)[]} (\texttt{48001}, \texttt{48002}), \texttt{401}, and \texttt{500} with \texttt{(DATABASE\_ERROR | UNKNOWN\_ERROR)[]}. \texttt{studitopia.vote\_forum\_content(puser\_id: uuid, ptarget\_type: string, ptarget\_id: uuid, pvote: integer)} returns \texttt{\{vote\_score: integer\}} and declares \texttt{UNKNOWN\_ERROR: 45000}, \texttt{POST\_NOT\_FOUND: 48001}, \texttt{COMMENT\_NOT\_FOUND: 48002}, \texttt{INVALID\_VOTE: 48005}, \texttt{DUPLICATED\_VOTE: 48006}, and \texttt{FORBIDDEN: 48010}.
\item \textbf{Vote state and notifications:} \texttt{list\_post.md} recommends \texttt{user\_vote} as the current user's \texttt{-1 | 0 | 1} state; \texttt{post\_management.md} says change/removal uses the same endpoint with \texttt{vote: 0}. \texttt{FR-38.1} and \texttt{NFR-PERF-6} include votes in forum activity notifications, but no vote-specific notification endpoint is defined in \texttt{API.md}.
\end{itemize}
\par\medskip

\hypertarget{sec:b11-source-alignment-notes}{}
\addcontentsline{toc}{subsubsection}{B11. Source Alignment Notes}
\subsubsection*{B11. Source Alignment Notes}
\par\medskip

\begin{itemize}
\item UC-04 lists voting as a downstream action after a Question or General Discussion post is published; UC-17 covers the vote action itself and does not repeat post-creation steps. UC-05 covers comment and answer creation and records vote context as related; UC-17 votes on those comments and answers without adding comments.
\item \texttt{AC-19.3} prevents more than one vote per target, \texttt{AC-19.4} allows changing or removing a vote, and \texttt{API.md} exposes both a \texttt{409 DUPLICATED\_VOTE} contract and \texttt{vote: 0} removal; \texttt{post\_management.md} resolves this at the UI layer by reflecting current state and offering change/removal.
\item \texttt{functional\_requirement.md} and \texttt{list\_post.md} refer to ranking updates, but \texttt{API.md} returns only \texttt{vote\_score} from the vote endpoint; no ranking response field is specified. \texttt{FR-38.1} includes votes in forum notifications, but no notification payload or endpoint is specified in the read sources.
\item Route names and component names in \texttt{list\_post.md} and \texttt{post\_management.md} are recommendations; \texttt{API.md} defines the exact endpoint and codes used above. No distinct included, extending, or specialized use case is named for US-19, so no new I-xx/E-xx/G-xx ID is introduced.
\end{itemize}
\par\medskip

\par\medskip
\hypertarget{sec:ui-prototype-specimens}{}
\addcontentsline{toc}{subsection}{UI Prototype Specimens}
\subsection*{UI Prototype Specimens}
\par\medskip
\begin{figure}[H]
\centering
\includegraphics[width=0.75\textwidth,height=0.42\textheight,keepaspectratio]{"../../prototype/forum_feed_interaction_states_desktop.png"}
\par\vspace{2pt}
{\footnotesize \textbf{UI Prototype:} Forum Feed - Interaction states - Desktop}
\end{figure}
\par\smallskip

\end{document}
